% FIXME comparison!
\subsection{メモリ「スナップショット」の比較}
\label{snapshots_comparing}

変化を見るために2つのメモリスナップショットを直接比較する技術は、8ビットコンピューターゲームをハックしたり、「ハイスコア」ファイルをハックしたりするためによく使用されました。

例えば、8ビットコンピューター（これらにはあまりメモリがありませんが、ゲームは通常さらに少ないメモリしか消費しません）にゲームがロードされていて、現在100発の弾があることがわかっている場合、すべてのメモリの「スナップショット」を取り、どこかにバックアップできます。次に一度撃つと、弾数が99になり、2番目の「スナップショット」を取って両方を比較します：最初に100だったバイトがどこかにあり、現在は99になっているはずです。

これらの8ビットゲームはしばしばアセンブリ言語で書かれ、そのような変数はグローバルだったという事実を考慮すると、メモリ内のどのアドレスが弾数を保持しているかを確実に言うことができます。逆アセンブルされたゲームコード内でそのアドレスへのすべての参照を検索すると、弾数を\glslink{decrement}{デクリメント}するコード片を見つけるのはそれほど難しくありませんでした。そして、そこに\gls{NOP}命令、または数個の\gls{NOP}を書き込み、永遠に100発の弾を持つゲームを手に入れることができました。
\myindex{BASIC!POKE}
これらの8ビットコンピューターのゲームは一般的に一定のアドレスにロードされ、また、各ゲームの異なるバージョンはあまりありませんでした（一般的に長期間1つのバージョンだけが人気でした）。そのため、熱心なゲーマーは、ハックするためにどのアドレスでどのバイトを上書きする必要があるか（BASICの\gls{POKE}命令を使用）を知っていました。これにより、8ビットゲーム関連の雑誌で公開された\gls{POKE}命令を含む「チート」リストが生まれました。参照：\href{http://go.yurichev.com/17114}{wikipedia}。

\myindex{MS-DOS}

同様に、「ハイスコア」ファイルを変更するのは簡単で、これは8ビットゲームだけではありません。スコア数に注目し、ファイルをどこかにバックアップしてください。「ハイスコア」数が変わったら、2つのファイルを比較するだけです。DOSユーティリティFC\footnote{バイナリファイルを比較するMS-DOSユーティリティ}（「ハイスコア」ファイルはしばしばバイナリ形式）でも実行できます。

いくつかのバイトが異なる点があり、どれがスコア番号を保持しているかを見るのは簡単です。
ただし、ゲーム開発者はそのようなトリックを完全に認識しており、プログラムをそれに対して保護する可能性があります。

この本の類似例：\myref{Millenium_DOS_game}。

% TODO: пример с какой-то простой игрушкой?

\subsubsection{Windowsレジストリ}

プログラムインストールの前後でWindowsレジストリを比較することも可能です。

これは、プログラムで使用されるレジストリ要素を見つける非常に人気のある方法です。
おそらく、これが「windowsレジストリクリーナー」シェアウェアが非常に人気である理由です。

\subsubsection{ブリンクコンパレーター}

ファイルやメモリスナップショットの比較は、ブリンクコンパレーター
\footnote{\url{http://go.yurichev.com/17348}}を思い出させます：
過去に天文学者によって使用された、移動する天体を見つけるための装置です。

ブリンクコンパレーターは、異なる時間に撮影された2つの写真を素早く切り替えることができ、
天文学者は視覚的に違いを見つけることができました。

ちなみに、冥王星は1930年にブリンクコンパレーターによって発見されました。